
\documentclass{article}
\usepackage{colortbl}
\usepackage{makecell}
\usepackage{multirow}
\usepackage{supertabular}

\begin{document}

\newcounter{utterance}

\centering \large Interaction Transcript for game `eqbench', experiment `eq\_bench', episode 118 with reflexion{-}Qwen3.5{-}9B.
\vspace{24pt}

{ \footnotesize  \setcounter{utterance}{1}
\setlength{\tabcolsep}{0pt}
\begin{supertabular}{c@{$\;$}|p{.15\linewidth}@{}p{.15\linewidth}p{.15\linewidth}p{.15\linewidth}p{.15\linewidth}p{.15\linewidth}}
    \# & \multicolumn{2}{c}{Player} && \multicolumn{2}{c}{Game Master} \\
    \hline

    \theutterance \stepcounter{utterance}  
    & & & \multicolumn{4}{p{0.6\linewidth}}{
        \cellcolor[rgb]{0.9,0.9,0.9}{
            \makecell[{{p{\linewidth}}}]{
                \texttt{\tiny{[P1$\langle$GM]}}
                \\ 
\texttt{Your task is to predict the likely emotional responses of a character in this dialogue:} \\
\\ 
\texttt{Freddy: You've got a fire in you, Izzy. Reminds me of a young me.} \\
\texttt{Izzy: Save your flattery, Lawson. It won't work on me.} \\
\texttt{Freddy: Who said anything about flattery? I'm merely stating an observation.} \\
\texttt{Izzy: Well, here's an observation for you. I'm not you, and I never will be.} \\
\texttt{Freddy: Oh, I'm well aware. You see, I win my cases.} \\
\texttt{Izzy: That's because you bend the law to your will. You manipulate, you lie.} \\
\texttt{Freddy: I do what it takes to win. Something you could learn from.} \\
\texttt{Izzy: I'd rather lose than become a monster like you.} \\
\texttt{Freddy: A monster? That's a bit harsh, don't you think?} \\
\texttt{Izzy: Not harsh enough. You've lost sight of what justice is.} \\
\texttt{Freddy: Justice is a luxury we can't afford, Izzy.} \\
\texttt{Izzy: No, justice is a necessity you've forgotten.} \\
\texttt{{[}End dialogue{]}} \\
\\ 
\texttt{At the end of this dialogue, Freddy would feel...} \\
\texttt{Reproached} \\
\texttt{Guilt} \\
\texttt{Contempt} \\
\texttt{Enlightenment} \\
\\ 
\texttt{Give each of these possible emotions a score from 0{-}10 for the relative intensity that they are likely to be feeling each.} \\
\\ 
\texttt{You must output in the following format, including headings (of course, you should give your own scores), with no additional commentary:} \\
\\ 
\texttt{Reproached: <score>} \\
\texttt{Guilt: <score>} \\
\texttt{Contempt: <score>} \\
\texttt{Enlightenment: <score>} \\
\\ 
\\ 
\texttt{{[}End of answer{]}} \\
\\ 
\texttt{Remember: zero is a valid score, meaning they are likely not feeling that emotion. You must score at least one emotion > 0.} \\
\\ 
\texttt{Your answer:} \\
            }
        }
    }
     \\ \\

    \theutterance \stepcounter{utterance}  
    & \multicolumn{4}{p{0.6\linewidth}}{
        \cellcolor[rgb]{0.561,0.737,0.561}{
            \makecell[{{p{\linewidth}}}]{
                \texttt{\tiny{[P1$\rangle$GM]}}
                \texttt{Reproached: 2} \\
\texttt{Guilt: 1} \\
\texttt{Contempt: 8} \\
\texttt{Enlightenment: 0} \\
            }
        }
    }
    & & \\ \\

    \theutterance \stepcounter{utterance}  
    & & & \multicolumn{2}{p{0.3\linewidth}}{
        \cellcolor[rgb]{0.9,0.9,0.9}{
            \makecell[{{p{\linewidth}}}]{
                \texttt{\tiny{[GM$|$GM]}}
                \texttt{Reproached: 2} \\
\texttt{Guilt: 1} \\
\texttt{Contempt: 8} \\
\texttt{Enlightenment: 0} \\
            }
        }
    }
    & & \\ \\

    \theutterance \stepcounter{utterance}  
    & & & \multicolumn{2}{p{0.3\linewidth}}{
        \cellcolor[rgb]{0.9,0.9,0.9}{
            \makecell[{{p{\linewidth}}}]{
                \texttt{\tiny{[GM$|$GM]}}
                \texttt{Reproached: 8} \\
\texttt{Guilt: 0} \\
\texttt{Contempt: 7} \\
\texttt{Enlightenment: 0} \\
            }
        }
    }
    & & \\ \\

    \theutterance \stepcounter{utterance}  
    & & & \multicolumn{2}{p{0.3\linewidth}}{
        \cellcolor[rgb]{0.9,0.9,0.9}{
            \makecell[{{p{\linewidth}}}]{
                \texttt{\tiny{[GM$|$GM]}}
                \texttt{game\_result = LOSE} \\
            }
        }
    }
    & & \\ \\

    \theutterance \stepcounter{utterance}  
    & & \multicolumn{4}{p{0.6\linewidth}}{
        \cellcolor[rgb]{0.95,0.95,0.95}{
            \makecell[{{p{\linewidth}}}]{
                \texttt{\tiny{[P1$\leftrightarrow$P]}}
                \texttt{Post{-}episode Reflexion (not sent to GM):} \\
\texttt{New plan: I failed because I scored "Contempt" too high (8) and "Reproached" too low (2), missing the nuance of Freddy's defensive yet self{-}justifying tone. Freddy is not genuinely feeling contempt for Izzy; he is using it as a rhetorical shield to deflect her moral accusations. His primary emotional state is a defensive reaction to being called a liar and a monster, which aligns more with feeling reproached (challenged in his integrity) or perhaps a mix of defensiveness and arrogance. I should have analyzed the subtext of "A monster? That's a bit harsh" as a deflection rather than a genuine feeling of superiority. Next time, I will lower the score for Contempt to reflect his performative nature and increase the score for Reproached to capture his defensive stance against Izzy's direct moral challenge, ensuring the highest score reflects the most immediate, reactive emotion triggered by the dialogue's climax.} \\
            }
        }
    }
    & \\ \\

\end{supertabular}
}

\end{document}
